\documentclass[11pt]{article}
\usepackage{amsmath, amssymb, geometry, enumitem, booktabs, hyperref}
\geometry{letterpaper, margin=1in}

\title{Linear Algebra: Comprehensive Curated Problem Sets \\ \large Based on Systems of Linear Equations \& Row Reduction}
\author{Math Curriculum Generator}
\date{}

\begin{document}

\maketitle
\tableofcontents
\newpage

% ==========================================
% CONCEPT CHECKLIST & QUESTION MAPPING
% ==========================================
\section{Concept Checklist \& Question Mapping}

\subsection*{Section 1.1: Systems of Linear Equations}
\begin{table}[h]
\centering
\begin{tabular}{@{}ll@{}}
\toprule
\textbf{Concept / Problem Type} & \textbf{Problem Numbers} \\ \midrule
C1.1: Direct Solution of Fully Determined Systems & 1.1.1 -- 1.1.6 \\
C1.2: Interpreting Reduced Matrices \& Consistency & 1.1.7 -- 1.1.10 \\
C1.3: Consistency Parameters Analysis ($h, k$, etc.) & 1.1.11 -- 1.1.14 \\
C1.4: Systems with Infinitely Many Solutions & 1.1.15 -- 1.1.20 \\
C1.5: Theoretical \& Conceptual Framework & 1.1.21 -- 1.1.24 \\
C1.6: Real-World Application and Modeling & 1.1.25 -- 1.1.27 \\
C1.7: Advanced Diagnostic Testing ("Find the Flaw") & 1.1.28 -- 1.1.30 \\ \bottomrule
\end{tabular}
\end{table}

\subsection*{Section 1.2: Row Reduction and Echelon Forms}
\begin{table}[h]
\centering
\begin{tabular}{@{}ll@{}}
\toprule
\textbf{Concept / Problem Type} & \textbf{Problem Numbers} \\ \midrule
C2.1: Identifying Echelon Forms (REF vs RREF) & 1.2.1 -- 1.2.5 \\
C2.2: Row Reduction Algorithm \& Pivot Identification & 1.2.6 -- 1.2.10 \\
C2.3: Extracting General Parametric Solutions & 1.2.11 -- 1.2.16 \\
C2.4: Parameter Analysis for System Behavior & 1.2.17 -- 1.2.21 \\
C2.5: Theoretical Implications of Pivots & 1.2.22 -- 1.2.25 \\
C2.6: Real-World Modeling (Traffic, Interpolation) & 1.2.26 -- 1.2.27 \\
C2.7: Advanced Diagnostic Testing ("Find the Flaw") & 1.2.28 -- 1.2.30 \\ \bottomrule
\end{tabular}
\end{table}

\newpage

% ==========================================
% PROBLEM SETS
% ==========================================
\section{Problem Sets}

\subsection{Homework 1.1: Systems of Linear Equations}

\textbf{Direct Solution of Fully Determined Systems}
\begin{enumerate}[label=\textbf{1.1.\arabic*}]
    \item Solve the linear system: 
    $x_1 + 2x_2 = 5$, 
    $2x_1 + 3x_2 = 8$.
    \item Solve the linear system using elementary row operations: 
    $3x_1 - x_2 = 7$, 
    $2x_1 + 4x_2 = 14$.
    \item Solve the $3 \times 3$ system:
    $x_1 + x_2 + x_3 = 6$,
    $2x_1 - x_2 + 3x_3 = 9$,
    $x_1 + 2x_2 - x_3 = 2$.
    \item Solve the system (Hint: A row swap may be required):
    $2x_2 + x_3 = 7$,
    $x_1 - x_2 + x_3 = 2$,
    $3x_1 + 2x_2 - x_3 = 4$.
    \item Solve the system by utilizing row scaling to avoid fractions:
    $2x_1 + 4x_2 - 2x_3 = -8$,
    $3x_1 - x_2 + x_3 = 11$,
    $x_1 + x_2 + x_3 = 5$.
    \item Solve the system corresponding to the given augmented matrix:
    $\left[\begin{array}{ccc|c} 1 & 2 & 3 & 14 \\ 0 & 2 & -1 & 1 \\ 0 & 0 & 4 & 12 \end{array}\right]$.
\end{enumerate}

\textbf{Interpreting Reduced Matrices \& Consistency}
\begin{enumerate}[label=\textbf{1.1.\arabic*}, resume]
    \item The augmented matrix of a linear system has been reduced to $\left[\begin{array}{ccc|c} 1 & 2 & 3 & 4 \\ 0 & 1 & 2 & 3 \\ 0 & 0 & 0 & 5 \end{array}\right]$. Describe the solution set.
    \item Describe the solution set of the system represented by $\left[\begin{array}{ccc|c} 1 & 0 & 0 & 2 \\ 0 & 1 & 0 & -1 \\ 0 & 0 & 1 & 4 \end{array}\right]$.
    \item A system reduces to $\left[\begin{array}{ccc|c} 1 & -1 & 2 & 4 \\ 0 & 1 & -1 & 2 \\ 0 & 0 & 0 & 0 \end{array}\right]$. Is the system consistent? Does it have a unique solution?
    \item Determine if the following system is consistent without completely solving it:
    $x_1 - x_2 = 1$, 
    $-2x_1 + 2x_2 = -2$, 
    $3x_1 - 3x_2 = 4$.
\end{enumerate}

\textbf{Consistency Parameters Analysis}
\begin{enumerate}[label=\textbf{1.1.\arabic*}, resume]
    \item Determine the value(s) of $h$ such that the augmented matrix represents a consistent linear system: $\left[\begin{array}{cc|c} 1 & 2 & 5 \\ 3 & h & 15 \end{array}\right]$.
    \item Determine the value(s) of $h$ such that the matrix represents a consistent system: $\left[\begin{array}{cc|c} 1 & 2 & 5 \\ 3 & h & 10 \end{array}\right]$.
    \item Choose $h$ and $k$ such that the system is consistent: $\left[\begin{array}{cc|c} 1 & 3 & h \\ -2 & -6 & k \end{array}\right]$.
    \item Find an equation involving $g, h$, and $k$ that makes this augmented matrix correspond to a consistent system: $\left[\begin{array}{ccc|c} 1 & -1 & 2 & g \\ 0 & 2 & 1 & h \\ 1 & 1 & 3 & k \end{array}\right]$.
\end{enumerate}

\textbf{Systems with Infinitely Many Solutions}
\begin{enumerate}[label=\textbf{1.1.\arabic*}, resume]
    \item Find the general parametric solution to the system:
    $x_1 - 2x_2 + x_3 = 0$, 
    $2x_1 - 3x_2 + x_3 = 2$, 
    $3x_1 - 5x_2 + 2x_3 = 2$.
    \item Find the general parametric solution to: $x_1 + 3x_2 - x_3 = 4$, $2x_1 + 6x_2 - 2x_3 = 8$.
    \item Express the solution to the single equation system parametrically: $x_1 + x_2 = 5$.
    \item A $3 \times 4$ augmented matrix reduces to $\left[\begin{array}{ccc|c} 1 & 2 & -1 & 4 \\ 0 & 0 & 0 & 0 \\ 0 & 0 & 0 & 0 \end{array}\right]$. Write the general solution.
    \item Find the general solution to: $x_1 + x_2 + x_3 + x_4 = 10$, $2x_1 + 2x_2 + 2x_3 + 2x_4 = 20$.
    \item Solve the system: $\left[\begin{array}{cccc|c} 1 & 0 & 2 & -1 & 4 \\ 0 & 1 & -1 & 3 & 5 \\ 0 & 0 & 0 & 0 & 0 \end{array}\right]$.
\end{enumerate}

\textbf{Theoretical \& Conceptual Framework}
\begin{enumerate}[label=\textbf{1.1.\arabic*}, resume]
    \item True or False: Multiplying a row by 0 is a valid elementary row operation. Justify.
    \item True or False: If two augmented matrices are row equivalent, they represent linear systems with the exact same solution set.
    \item Geometrically, what does the solution to a system of three linear equations in three variables represent if the system yields exactly one solution?
    \item True or False: The elementary row operation $R_1 \leftarrow R_1 + R_2$ followed by $R_2 \leftarrow R_2 + R_1$ leaves the matrix entirely unchanged.
\end{enumerate}

\textbf{Real-World Application and Modeling}
\begin{enumerate}[label=\textbf{1.1.\arabic*}, resume]
    \item \textbf{Arbitrage Pricing Setup:} Asset A pays $\$10$ in a Bull market and $\$5$ in a Bear market. Its current price is $\$8$. Asset B pays $\$20$ in both markets. Its current price is $\$18$. Let $p_1, p_2$ be the state prices of the Bull and Bear markets respectively. Write the augmented matrix for this system.
    \item \textbf{Arbitrage Pricing Solve:} Solve the matrix from Problem 1.1.25 to find the risk-neutral state prices $p_1$ and $p_2$.
    \item \textbf{Geometric Intersection:} Find the point $(x_1, x_2)$ that lies on the line $3x_1 + 4x_2 = 10$ and on the line $x_1 - 2x_2 = 0$.
\end{enumerate}

\textbf{Advanced Diagnostic Testing: "Find the Flaw"}
\begin{enumerate}[label=\textbf{1.1.\arabic*}, resume]
    \item \textbf{Flawed Solution:} Solve $x_1+2x_2=4, 3x_1+4x_2=10$. \textit{Student work:} Matrix is $\left[\begin{array}{cc|c} 1 & 2 & 4 \\ 3 & 4 & 10 \end{array}\right]$. Apply $R_2 \leftarrow R_2 - 3R_1$. The new $R_2$ is $[0, -2 | 22]$. So $-2x_2 = 22 \implies x_2 = -11$. Find the flaw and provide the correct solution.
    \item \textbf{Flawed Solution:} Matrix $\left[\begin{array}{cc|c} 1 & 2 & 3 \\ 0 & 0 & 0 \end{array}\right]$. \textit{Student work:} Because the bottom row is $0=0$, the variables vanish, meaning there is no solution. Find the flaw and correct it.
    \item \textbf{Flawed Solution:} Matrix $\left[\begin{array}{cc|c} 1 & -1 & 2 \\ 2 & -2 & 5 \end{array}\right]$. \textit{Student work:} There are no zero rows, and no row of the form $[0, 0 | b]$, so the system is consistent and has a unique solution. Find the flaw and correct it.
\end{enumerate}

% ---------------------------------------------------------
\subsection{Homework 1.2: Row Reduction and Echelon Forms}

\textbf{Identifying Echelon Forms}
\begin{enumerate}[label=\textbf{1.2.\arabic*}]
    \item Classify the matrix as in Reduced Echelon Form (RREF), Echelon Form only (REF), or neither: $\left[\begin{array}{ccc} 1 & 2 & 0 \\ 0 & 0 & 1 \\ 0 & 0 & 0 \end{array}\right]$.
    \item Classify the matrix: $\left[\begin{array}{ccc} 1 & 0 & 3 \\ 0 & 2 & 0 \\ 0 & 0 & 0 \end{array}\right]$.
    \item Classify the matrix: $\left[\begin{array}{ccc} 0 & 1 & 0 \\ 1 & 0 & 0 \\ 0 & 0 & 1 \end{array}\right]$.
    \item Classify the matrix: $\left[\begin{array}{cc} 1 & 2 \\ 0 & 1 \end{array}\right]$.
    \item Classify the matrix: $\left[\begin{array}{ccc} 1 & 0 & 0 \\ 0 & 1 & 0 \\ 0 & 0 & 1 \end{array}\right]$.
\end{enumerate}

\textbf{Row Reduction Algorithm \& Pivot Identification}
\begin{enumerate}[label=\textbf{1.2.\arabic*}, resume]
    \item Row reduce to RREF and identify the pivot columns: $\left[\begin{array}{cc} 1 & 3 \\ 2 & 6 \end{array}\right]$.
    \item Row reduce to RREF and identify the pivot positions: $\left[\begin{array}{ccc} 1 & 2 & 3 \\ 4 & 5 & 6 \end{array}\right]$.
    \item Row reduce to RREF: $\left[\begin{array}{ccc} 0 & 2 & 4 \\ 1 & 3 & 5 \end{array}\right]$.
    \item Row reduce to RREF: $\left[\begin{array}{ccc} 1 & -1 & 2 \\ -2 & 2 & -4 \end{array}\right]$.
    \item A matrix is in REF: $\left[\begin{array}{ccc|c} 1 & 2 & -1 & 4 \\ 0 & 1 & 3 & 5 \\ 0 & 0 & 2 & 6 \end{array}\right]$. Complete the backward phase to find the RREF.
\end{enumerate}

\textbf{Extracting General Parametric Solutions}
\begin{enumerate}[label=\textbf{1.2.\arabic*}, resume]
    \item Find the general solution of the system whose augmented matrix is $\left[\begin{array}{ccc|c} 1 & 0 & 4 & 2 \\ 0 & 1 & -2 & 3 \end{array}\right]$.
    \item Find the general solution: $\left[\begin{array}{ccc|c} 1 & 0 & 0 & 5 \\ 0 & 0 & 1 & 4 \end{array}\right]$.
    \item Find the general solution: $\left[\begin{array}{ccc|c} 1 & 2 & 0 & 0 \\ 0 & 0 & 1 & 0 \\ 0 & 0 & 0 & 1 \end{array}\right]$.
    \item Find the general solution: $\left[\begin{array}{ccc|c} 0 & 1 & 0 & 2 \\ 0 & 0 & 1 & 3 \end{array}\right]$.
    \item The matrix is in REF. Find the general solution: $\left[\begin{array}{ccc|c} 1 & 2 & 3 & 4 \\ 0 & 1 & 2 & 3 \end{array}\right]$.
    \item The matrix is in REF. Find the general solution: $\left[\begin{array}{ccc|c} 2 & 4 & 2 & 6 \\ 0 & 0 & 3 & 9 \end{array}\right]$.
\end{enumerate}

\textbf{Parameter Analysis for System Behavior}
\begin{enumerate}[label=\textbf{1.2.\arabic*}, resume]
    \item Determine the value(s) of $h$ such that the system $\left[\begin{array}{cc|c} 1 & -2 & 3 \\ 0 & h-1 & 4 \end{array}\right]$ has no solution.
    \item Determine $h$ and $k$ such that the system $\left[\begin{array}{cc|c} 1 & 3 & 2 \\ 2 & h & k \end{array}\right]$ has infinitely many solutions.
    \item Determine $h$ and $k$ such that the system from 1.2.18 has a unique solution.
    \item Determine $h$ such that $\left[\begin{array}{cc|c} 1 & h & 2 \\ -2 & -4 & -4 \end{array}\right]$ has infinitely many solutions.
    \item Consider the system $x_1 + hx_2 = 2$, $3x_1 + 6x_2 = k$. For what value of $h$ is it possible for the system to be inconsistent (depending on $k$)?
\end{enumerate}

\textbf{Theoretical Implications of Pivots}
\begin{enumerate}[label=\textbf{1.2.\arabic*}, resume]
    \item Suppose a $4 \times 6$ augmented matrix has 4 pivots in its coefficient matrix portion. Is the system consistent?
    \item True or False: A matrix can be row reduced to multiple different echelon forms, but only one reduced echelon form.
    \item If a coefficient matrix has a pivot in every column, what can be said about the number of free variables?
    \item Suppose a system of 3 equations in 5 variables is consistent. What is the minimum number of free variables the system can have?
\end{enumerate}

\textbf{Real-World Modeling}
\begin{enumerate}[label=\textbf{1.2.\arabic*}, resume]
    \item \textbf{Traffic Flow:} Four intersections (A, B, C, D) form a square one-way loop. 
    Node A: In 200, Out $x_1 + x_4$.
    Node B: In $x_1 + 50$, Out $x_2$.
    Node C: In $x_2$, Out $x_3 + 150$.
    Node D: In $x_3 + x_4$, Out 100.
    Set up the augmented matrix and find the general solution for the traffic flow $x_1, x_2, x_3, x_4$.
    \item \textbf{Polynomial Interpolation:} Find the polynomial of degree 2, $p(t) = a_0 + a_1 t + a_2 t^2$, whose graph passes through $(1, 4)$, $(2, 9)$, and $(3, 16)$. Set up the matrix and solve for $a_0, a_1, a_2$.
\end{enumerate}

\textbf{Advanced Diagnostic Testing: "Find the Flaw"}
\begin{enumerate}[label=\textbf{1.2.\arabic*}, resume]
    \item \textbf{Flawed Solution:} Matrix $\left[\begin{array}{cc|c} 1 & 2 & 3 \\ 0 & 1 & 2 \end{array}\right]$. \textit{Student work:} To do the backward phase, I need to clear the 2. $R_2 \leftarrow R_2 - 0.5 R_1$ yields $[-0.5, 0 | 0.5]$. Find the flaw and correctly find the RREF.
    \item \textbf{Flawed Solution:} Matrix $\left[\begin{array}{ccc|c} 1 & 0 & 3 & 4 \\ 0 & 0 & 0 & 0 \end{array}\right]$. \textit{Student work:} Column 2 is all zeros, which means $x_2 = 0$. Find the flaw and correct it.
    \item \textbf{Flawed Solution:} Matrix $\left[\begin{array}{ccc|c} 1 & 0 & -2 & 5 \\ 0 & 1 & 1 & 2 \end{array}\right]$. \textit{Student work:} $x_1 - 2x_3 = 5 \implies x_3 = \frac{1}{2}x_1 - \frac{5}{2}$. Find the flaw and correct it.
\end{enumerate}

\newpage
% ==========================================
% DETAILED SOLUTIONS
% ==========================================
\section{Detailed Solutions}

\subsection{Solutions to Homework 1.1}
\begin{enumerate}[label=\textbf{Sol 1.1.\arabic*}]
    \item $\left[\begin{array}{cc|c} 1 & 2 & 5 \\ 2 & 3 & 8 \end{array}\right] \xrightarrow{R_2 - 2R_1} \left[\begin{array}{cc|c} 1 & 2 & 5 \\ 0 & -1 & -2 \end{array}\right]$. $x_2 = 2$. $x_1 + 2(2) = 5 \implies x_1 = 1$. \textbf{Answer:} $(1, 2)$.
    \item $\left[\begin{array}{cc|c} 3 & -1 & 7 \\ 2 & 4 & 14 \end{array}\right] \xrightarrow{\frac{1}{2}R_2} \left[\begin{array}{cc|c} 3 & -1 & 7 \\ 1 & 2 & 7 \end{array}\right] \xrightarrow{R_1 \leftrightarrow R_2} \left[\begin{array}{cc|c} 1 & 2 & 7 \\ 3 & -1 & 7 \end{array}\right] \xrightarrow{R_2 - 3R_1} \left[\begin{array}{cc|c} 1 & 2 & 7 \\ 0 & -7 & -14 \end{array}\right]$. $x_2 = 2$. $x_1 + 4 = 7 \implies x_1 = 3$. \textbf{Answer:} $(3, 2)$.
    \item $\left[\begin{array}{ccc|c} 1 & 1 & 1 & 6 \\ 2 & -1 & 3 & 9 \\ 1 & 2 & -1 & 2 \end{array}\right]$. $R_2-2R_1$, $R_3-R_1 \implies \left[\begin{array}{ccc|c} 1 & 1 & 1 & 6 \\ 0 & -3 & 1 & -3 \\ 0 & 1 & -2 & -4 \end{array}\right]$. Swap $R_2, R_3 \implies \left[\begin{array}{ccc|c} 1 & 1 & 1 & 6 \\ 0 & 1 & -2 & -4 \\ 0 & -3 & 1 & -3 \end{array}\right]$. $R_3+3R_2 \implies \left[\begin{array}{ccc|c} 1 & 1 & 1 & 6 \\ 0 & 1 & -2 & -4 \\ 0 & 0 & -5 & -15 \end{array}\right]$. $x_3 = 3$. $x_2 - 6 = -4 \implies x_2 = 2$. $x_1 + 2 + 3 = 6 \implies x_1 = 1$. \textbf{Answer:} $(1, 2, 3)$.
    \item $\left[\begin{array}{ccc|c} 0 & 2 & 1 & 7 \\ 1 & -1 & 1 & 2 \\ 3 & 2 & -1 & 4 \end{array}\right]$. Swap $R_1, R_2$, then $R_3-3R_1 \implies \left[\begin{array}{ccc|c} 1 & -1 & 1 & 2 \\ 0 & 2 & 1 & 7 \\ 0 & 5 & -4 & -2 \end{array}\right]$. $R_3 \leftarrow 2R_3 - 5R_2 \implies \left[\begin{array}{ccc|c} 1 & -1 & 1 & 2 \\ 0 & 2 & 1 & 7 \\ 0 & 0 & -13 & -39 \end{array}\right]$. $x_3=3$. $2x_2+3=7 \implies x_2=2$. $x_1-2+3=2 \implies x_1=1$. \textbf{Answer:} $(1, 2, 3)$.
    \item Scale $R_1$ by $1/2 \implies x_1 + 2x_2 - x_3 = -4$. Matrix: $\left[\begin{array}{ccc|c} 1 & 2 & -1 & -4 \\ 3 & -1 & 1 & 11 \\ 1 & 1 & 1 & 5 \end{array}\right]$. $R_2-3R_1, R_3-R_1 \implies \left[\begin{array}{ccc|c} 1 & 2 & -1 & -4 \\ 0 & -7 & 4 & 23 \\ 0 & -1 & 2 & 9 \end{array}\right]$. Swap $R_2, R_3$ and negate $\implies \left[\begin{array}{ccc|c} 1 & 2 & -1 & -4 \\ 0 & 1 & -2 & -9 \\ 0 & -7 & 4 & 23 \end{array}\right]$. $R_3+7R_2 \implies -10x_3 = -40 \implies x_3=4$. $x_2 - 8 = -9 \implies x_2 = -1$. $x_1 - 2 - 4 = -4 \implies x_1 = 2$. \textbf{Answer:} $(2, -1, 4)$.
    \item Matrix is already in echelon form. $R_3 \implies 4x_3 = 12 \implies x_3=3$. $R_2 \implies 2x_2 - 3 = 1 \implies x_2 = 2$. $R_1 \implies x_1 + 2(2) + 3(3) = 14 \implies x_1 = 1$. \textbf{Answer:} $(1, 2, 3)$.
    \item The bottom row is $0x_1 + 0x_2 + 0x_3 = 5 \implies 0 = 5$, a contradiction. \textbf{Answer:} The system is inconsistent (empty solution set).
    \item The matrix is reduced to the identity on the left. \textbf{Answer:} The system is consistent and has a unique solution $(2, -1, 4)$.
    \item Row 3 is $0=0$. This is universally true. Columns 1 and 2 are pivot columns. Column 3 lacks a pivot, meaning $x_3$ is a free variable. \textbf{Answer:} Consistent, with infinitely many solutions.
    \item $\left[\begin{array}{cc|c} 1 & -1 & 1 \\ -2 & 2 & -2 \\ 3 & -3 & 4 \end{array}\right]$. $R_2+2R_1 \implies [0, 0 | 0]$. $R_3-3R_1 \implies [0, 0 | 1]$. The last row $0=1$ is a contradiction. \textbf{Answer:} Inconsistent.
    \item $R_2 \leftarrow R_2 - 3R_1 \implies \left[\begin{array}{cc|c} 1 & 2 & 5 \\ 0 & h-6 & 0 \end{array}\right]$. The bottom row is $(h-6)x_2 = 0$. Whether $h=6$ (yielding $0=0$) or $h \neq 6$ (yielding $x_2=0$), there is no contradiction. \textbf{Answer:} Consistent for all real $h$.
    \item $R_2 \leftarrow R_2 - 3R_1 \implies \left[\begin{array}{cc|c} 1 & 2 & 5 \\ 0 & h-6 & -5 \end{array}\right]$. If $h-6=0$, we get $0=-5$ (inconsistent). \textbf{Answer:} Consistent for all $h \neq 6$.
    \item $R_2 \leftarrow R_2 + 2R_1 \implies \left[\begin{array}{cc|c} 1 & 3 & h \\ 0 & 0 & k+2h \end{array}\right]$. For consistency, we must have $0 = k + 2h$. \textbf{Answer:} $k = -2h$.
    \item $R_3 \leftarrow R_3 - R_1 \implies [0, 2, 1 | k-g]$. Notice this matches the left side of $R_2$. $R_3 \leftarrow R_3 - R_2 \implies [0, 0, 0 | (k-g) - h]$. For consistency, the right side must be 0. \textbf{Answer:} $k - g - h = 0$.
    \item Matrix reduces to $\left[\begin{array}{ccc|c} 1 & -2 & 1 & 0 \\ 0 & 1 & -1 & 2 \\ 0 & 0 & 0 & 0 \end{array}\right]$. $x_3$ is free. $x_2 = 2 + x_3$. $x_1 = 2x_2 - x_3 = 2(2+x_3) - x_3 = 4 + x_3$. \textbf{Answer:} $x_1 = 4 + x_3$, $x_2 = 2 + x_3$, $x_3$ is free.
    \item $\left[\begin{array}{ccc|c} 1 & 3 & -1 & 4 \\ 2 & 6 & -2 & 8 \end{array}\right] \xrightarrow{R_2-2R_1} \left[\begin{array}{ccc|c} 1 & 3 & -1 & 4 \\ 0 & 0 & 0 & 0 \end{array}\right]$. $x_2, x_3$ are free. \textbf{Answer:} $x_1 = 4 - 3x_2 + x_3$, $x_2$ free, $x_3$ free.
    \item $x_1 = 5 - x_2$. \textbf{Answer:} $x_1 = 5 - x_2$, $x_2$ is free.
    \item $x_1, x_2$ are basic. $x_3$ is free, $x_4$ is free. From R1: $x_1 + 2x_2 - x_3 = 4 \implies x_1 = 4 - 2x_2 + x_3$. \textbf{Answer:} $x_1 = 4 - 2x_2 + x_3$, $x_2$ free, $x_3$ free, $x_4$ free.
    \item Both equations represent the same plane. $x_1 = 10 - x_2 - x_3 - x_4$. \textbf{Answer:} $x_1 = 10 - x_2 - x_3 - x_4$, $x_2, x_3, x_4$ are free.
    \item Pivots in col 1, 2. Free vars $x_3, x_4$. $x_2 - x_3 + 3x_4 = 5 \implies x_2 = 5 + x_3 - 3x_4$. $x_1 + 2x_3 - x_4 = 4 \implies x_1 = 4 - 2x_3 + x_4$. \textbf{Answer:} Parametric solution listed.
    \item \textbf{False.} Scaling requires multiplication by a non-zero constant ($c \neq 0$). Multiplying by 0 wipes out information and cannot be reversed.
    \item \textbf{True.} Row equivalence is defined by sequences of operations that strictly preserve the underlying algebraic equality and therefore the solution set.
    \item \textbf{Answer:} A single, unique point in 3D coordinate space $(x_1, x_2, x_3)$ where all three planes intersect.
    \item \textbf{False.} The operations are applied sequentially. $R_1$ changes first, so the new $R_1$ is used to modify $R_2$, fundamentally altering the matrix.
    \item Variables are $p_1$ (Bull price), $p_2$ (Bear price). Asset A: $10p_1 + 5p_2 = 8$. Asset B: $20p_1 + 20p_2 = 18$. \textbf{Answer:} $\left[\begin{array}{cc|c} 10 & 5 & 8 \\ 20 & 20 & 18 \end{array}\right]$.
    \item $\left[\begin{array}{cc|c} 10 & 5 & 8 \\ 20 & 20 & 18 \end{array}\right] \xrightarrow{R_2-2R_1} \left[\begin{array}{cc|c} 10 & 5 & 8 \\ 0 & 10 & 2 \end{array}\right]$. $10p_2 = 2 \implies p_2 = 0.2$. $10p_1 + 5(0.2) = 8 \implies 10p_1 = 7 \implies p_1 = 0.7$. \textbf{Answer:} $p_1 = 0.7, p_2 = 0.2$.
    \item $\left[\begin{array}{cc|c} 3 & 4 & 10 \\ 1 & -2 & 0 \end{array}\right] \xrightarrow{R_1 \leftrightarrow R_2} \left[\begin{array}{cc|c} 1 & -2 & 0 \\ 3 & 4 & 10 \end{array}\right] \xrightarrow{R_2-3R_1} \left[\begin{array}{cc|c} 1 & -2 & 0 \\ 0 & 10 & 10 \end{array}\right]$. $x_2=1$. $x_1 - 2 = 0 \implies x_1 = 2$. \textbf{Answer:} $(2, 1)$.
    \item \textbf{Flaw:} The student botched the arithmetic of $R_2 - 3R_1$. They computed $10 - 3(4) = 10 - 12 = -2$, but wrote $22$ (adding instead of subtracting). \textbf{Correction:} The new $R_2$ is $[0, -2 | -2]$. $-2x_2 = -2 \implies x_2 = 1$. Then $x_1 + 2(1) = 4 \implies x_1 = 2$.
    \item \textbf{Flaw:} $0=0$ is a universal truth, not a contradiction. The system is consistent but underdetermined. \textbf{Correction:} $x_1 + 2x_2 = 3 \implies x_1 = 3 - 2x_2$, with $x_2$ as a free variable.
    \item \textbf{Flaw:} The student stopped without fully reducing. A matrix without zero rows can still be inconsistent. \textbf{Correction:} $R_2 - 2R_1 \implies [0, 0 | 1]$. This is a contradiction $0=1$. The system is actually inconsistent.
\end{enumerate}

\subsection{Solutions to Homework 1.2}
\begin{enumerate}[label=\textbf{Sol 1.2.\arabic*}]
    \item Leading 1s step to the right, zeros above and below leading 1s, zero row at the bottom. \textbf{Answer:} RREF.
    \item Leading entries step right, but leading entry in R2 is not 1 (it's 2). \textbf{Answer:} REF only.
    \item Nonzero rows are present, but leading entry of R2 is to the left of R1's leading entry. \textbf{Answer:} Neither.
    \item Leading entries are 1, but there is a non-zero number (2) above the leading 1 in column 2. \textbf{Answer:} REF only.
    \item Identity matrix satisfies all conditions perfectly. \textbf{Answer:} RREF.
    \item $\left[\begin{array}{cc} 1 & 3 \\ 2 & 6 \end{array}\right] \xrightarrow{R_2-2R_1} \left[\begin{array}{cc} 1 & 3 \\ 0 & 0 \end{array}\right]$. \textbf{Answer:} RREF is $\left[\begin{array}{cc} 1 & 3 \\ 0 & 0 \end{array}\right]$. Pivot column is 1.
    \item $\left[\begin{array}{ccc} 1 & 2 & 3 \\ 4 & 5 & 6 \end{array}\right] \xrightarrow{R_2-4R_1} \left[\begin{array}{ccc} 1 & 2 & 3 \\ 0 & -3 & -6 \end{array}\right] \xrightarrow{-\frac{1}{3}R_2} \left[\begin{array}{ccc} 1 & 2 & 3 \\ 0 & 1 & 2 \end{array}\right] \xrightarrow{R_1-2R_2} \left[\begin{array}{ccc} 1 & 0 & -1 \\ 0 & 1 & 2 \end{array}\right]$. \textbf{Answer:} RREF is $\left[\begin{array}{ccc} 1 & 0 & -1 \\ 0 & 1 & 2 \end{array}\right]$. Pivot positions are $(1,1)$ and $(2,2)$.
    \item Swap $R_1, R_2 \implies \left[\begin{array}{ccc} 1 & 3 & 5 \\ 0 & 2 & 4 \end{array}\right] \xrightarrow{\frac{1}{2}R_2} \left[\begin{array}{ccc} 1 & 3 & 5 \\ 0 & 1 & 2 \end{array}\right] \xrightarrow{R_1-3R_2} \left[\begin{array}{ccc} 1 & 0 & -1 \\ 0 & 1 & 2 \end{array}\right]$. \textbf{Answer:} Matrix provided.
    \item $R_2+2R_1 \implies \left[\begin{array}{ccc} 1 & -1 & 2 \\ 0 & 0 & 0 \end{array}\right]$. \textbf{Answer:} Matrix provided. Pivot in col 1.
    \item $R_2 \gets R_2/2 \implies [0, 0, 1 | 3]$. $R_2 \gets R_2 - 3R_3 \implies [0, 1, 0 | 5 - 9 = -4]$. $R_1 \gets R_1 + R_3 \implies [1, 2, 0 | 7]$. $R_1 \gets R_1 - 2R_2 \implies [1, 0, 0 | 7 - 2(-4) = 15]$. \textbf{Answer:} $\left[\begin{array}{ccc|c} 1 & 0 & 0 & 15 \\ 0 & 1 & 0 & -4 \\ 0 & 0 & 1 & 3 \end{array}\right]$.
    \item Basic $x_1, x_2$. Free $x_3$. $x_1 + 4x_3 = 2 \implies x_1 = 2 - 4x_3$. $x_2 - 2x_3 = 3 \implies x_2 = 3 + 2x_3$. \textbf{Answer:} Parametric equations listed.
    \item Basic $x_1, x_3$. Free $x_2$. \textbf{Answer:} $x_1 = 5$, $x_2$ is free, $x_3 = 4$.
    \item Row 3 is $0=1$, a contradiction. \textbf{Answer:} Inconsistent (no solution).
    \item Basic $x_2, x_3$. Free $x_1$. \textbf{Answer:} $x_1$ is free, $x_2 = 2, x_3 = 3$.
    \item Backward phase: $R_1 - 2R_2 \implies [1, 0, -1 | -2]$. Basic $x_1, x_2$. Free $x_3$. \textbf{Answer:} $x_1 = -2 + x_3, x_2 = 3 - 2x_3, x_3$ free.
    \item Scale $R_2$ by $1/3 \implies [0,0,1|3]$. $R_1 - 2R_2 \implies [2, 4, 0 | 0]$. Scale $R_1$ by $1/2 \implies [1, 2, 0 | 0]$. \textbf{Answer:} $x_1 = -2x_2, x_2$ free, $x_3 = 3$.
    \item Need a row $[0 \dots 0 | b \neq 0]$. Row 2 is $[0, h-1 | 4]$. This is a contradiction if $h-1 = 0 \implies h = 1$. \textbf{Answer:} $h = 1$.
    \item $R_2 - 2R_1 \implies [0, h-6 | k-4]$. For infinite solutions, we need no contradiction AND a free variable. So $h-6 = 0 \implies h=6$, and $k-4 = 0 \implies k=4$. \textbf{Answer:} $h = 6, k = 4$.
    \item For a unique solution, we need a pivot in every variable column. The coefficient of $x_2$ in $R_2$ must be non-zero. $h-6 \neq 0 \implies h \neq 6$. $k$ can be anything. \textbf{Answer:} $h \neq 6, k \in \mathbb{R}$.
    \item $R_2+2R_1 \implies [0, 2h-4 | 0]$. For infinite solutions, we need $2h-4 = 0 \implies h=2$. \textbf{Answer:} $h=2$.
    \item Matrix is $\left[\begin{array}{cc|c} 1 & h & 2 \\ 3 & 6 & k \end{array}\right]$. $R_2-3R_1 \implies [0, 6-3h | k-6]$. If $6-3h=0 \implies h=2$, the row becomes $[0, 0 | k-6]$. If $k \neq 6$, it's inconsistent. \textbf{Answer:} $h = 2$.
    \item \textbf{Yes.} A $4 \times 6$ augmented matrix has 4 rows and 5 coefficient columns. If the coefficient matrix has 4 pivots, every row has a pivot on the left side. The augmented column cannot have a pivot, so it's consistent.
    \item \textbf{True.} Uniqueness Theorem of Reduced Echelon Form guarantees only one RREF exists per matrix, while REF depends on choices made during reduction.
    \item \textbf{Answer:} There are zero free variables. Every variable is basic.
    \item \textbf{Answer:} 2 free variables. Max pivots is 3 (since there are 3 rows). 5 variables - 3 basic = 2 free minimum.
    \item $\left[\begin{array}{cccc|c} 1 & 0 & 0 & 1 & 200 \\ 1 & -1 & 0 & 0 & -50 \\ 0 & 1 & -1 & 0 & 150 \\ 0 & 0 & 1 & 1 & 100 \end{array}\right]$. $R_2-R_1 \implies [0, -1, 0, -1 | -250]$. $R_3+R_2 \implies [0, 0, -1, -1 | -100]$. $R_4+R_3 \implies [0, 0, 0, 0 | 0]$. \textbf{Answer:} $x_1 = 200-x_4$, $x_2 = 250-x_4$, $x_3 = 100-x_4$, $x_4$ free.
    \item $a_0+a_1+a_2=4$, $a_0+2a_1+4a_2=9$, $a_0+3a_1+9a_2=16$. RREF yields $a_0=1, a_1=2, a_2=1$. \textbf{Answer:} $p(t) = 1 + 2t + t^2$.
    \item \textbf{Flaw:} The student used $R_1$ to clear an entry in $R_2$. The backward phase requires using the pivot row to clear entries \textit{above} it. \textbf{Correction:} $R_1 \leftarrow R_1 - 2R_2 \implies [1, 0 | -1]$.
    \item \textbf{Flaw:} A column of zeros in the coefficient matrix does not mean the variable equals zero; it means the variable lacks a pivot and is therefore free. \textbf{Correction:} $x_1=4$, $x_3=0$, and $x_2$ is free.
    \item \textbf{Flaw:} The student expressed the free variable ($x_3$) in terms of the basic variable ($x_1$). Parametric solutions must express basic variables in terms of free ones. \textbf{Correction:} $x_1 = 5 + 2x_3$.
\end{enumerate}

\newpage
% ==========================================
% REPEATED CONCEPT CHECKLIST
% ==========================================
\section{Concept Checklist \& Question Mapping (Repeated)}

\subsection*{Section 1.1: Systems of Linear Equations}
\begin{table}[h]
\centering
\begin{tabular}{@{}ll@{}}
\toprule
\textbf{Concept / Problem Type} & \textbf{Problem Numbers} \\ \midrule
C1.1: Direct Solution of Fully Determined Systems & 1.1.1 -- 1.1.6 \\
C1.2: Interpreting Reduced Matrices \& Consistency & 1.1.7 -- 1.1.10 \\
C1.3: Consistency Parameters Analysis ($h, k$, etc.) & 1.1.11 -- 1.1.14 \\
C1.4: Systems with Infinitely Many Solutions & 1.1.15 -- 1.1.20 \\
C1.5: Theoretical \& Conceptual Framework & 1.1.21 -- 1.1.24 \\
C1.6: Real-World Application and Modeling & 1.1.25 -- 1.1.27 \\
C1.7: Advanced Diagnostic Testing ("Find the Flaw") & 1.1.28 -- 1.1.30 \\ \bottomrule
\end{tabular}
\end{table}

\subsection*{Section 1.2: Row Reduction and Echelon Forms}
\begin{table}[h]
\centering
\begin{tabular}{@{}ll@{}}
\toprule
\textbf{Concept / Problem Type} & \textbf{Problem Numbers} \\ \midrule
C2.1: Identifying Echelon Forms (REF vs RREF) & 1.2.1 -- 1.2.5 \\
C2.2: Row Reduction Algorithm \& Pivot Identification & 1.2.6 -- 1.2.10 \\
C2.3: Extracting General Parametric Solutions & 1.2.11 -- 1.2.16 \\
C2.4: Parameter Analysis for System Behavior & 1.2.17 -- 1.2.21 \\
C2.5: Theoretical Implications of Pivots & 1.2.22 -- 1.2.25 \\
C2.6: Real-World Modeling (Traffic, Interpolation) & 1.2.26 -- 1.2.27 \\
C2.7: Advanced Diagnostic Testing ("Find the Flaw") & 1.2.28 -- 1.2.30 \\ \bottomrule
\end{tabular}
\end{table}

\end{document}